\documentclass[11pt,a4paper]{article}

\usepackage[T1]{fontenc}
\usepackage[utf8]{inputenc}
\usepackage{lmodern}
\usepackage[english]{babel}
\usepackage[margin=2.5cm]{geometry}
\usepackage{booktabs}
\usepackage{tabularx}
\usepackage{graphicx}
\usepackage{float}
\usepackage{enumitem}
\usepackage{xcolor}
\usepackage{listings}
\usepackage[colorlinks=true,linkcolor=black,urlcolor=blue]{hyperref}

\renewcommand{\arraystretch}{1.25}

% Red reminder. Search for "TODO" before submitting.
\providecommand{\todo}[1]{\textcolor{red}{[TODO: #1]}}

\lstset{
  basicstyle=\ttfamily\footnotesize,
  breaklines=true,
  frame=single,
  columns=fullflexible
}

\begin{document}

% ------------------------------------------------------------------
% TITLE PAGE
% ------------------------------------------------------------------
\begin{titlepage}
  \centering
  {\large Universidad Distrital Francisco Jos\'e de Caldas\\
  School of Engineering -- Computer Engineering Program\par}
  \vspace{1.5cm}
  {\Large Data Analysis Programming -- Semester 2026-II\par}
  \vspace{1cm}
  {\LARGE\bfseries Workshop No.\ 1\\[0.3cm]
  Hazard Definition, Source Validation,\\and Data Contract\par}
  \vspace{1.5cm}
  {\Large Team 8 -- Tsunami Monitoring\par}
  \vspace{1cm}
  \begin{tabular}{ll}
    Lina Moreno & 20231020133 \\
    Jaider \'Avila & 20231020200 \\
    Gabriela Mart\'inez & 20231020205 \\
  \end{tabular}
  \vspace{1.5cm}

  Professor: Eng.\ Carlos Andr\'es Sierra, M.Sc.\par
  \vspace{0.5cm}
  Repository: \url{https://github.com/g8martinezsilva/tsunami-monitoring}\par
  \vfill
  \today
\end{titlepage}

\tableofcontents
\clearpage

% One file per section, so that each person edits a different file.

% ===== 00_summary.tex =====
% Owner: TBD
\section*{Executive Summary}
\todo{Five to eight lines: hazard, scope, the four sources, main findings of the profiling, and the proposed indicator.}



% ===== 01_hazard.tex =====
% Owner: TBD
\section{Hazard Definition}
\label{sec:hazard}

\subsection{Phenomenon and Severity Measure}
\todo{What a tsunami is, how it is caused, how it is measured. Severity measure of the system (earthquake magnitude and its type, depth in km) and meaning of the units.}

\subsection{Why Monitoring Matters}
\todo{Motivation. State that the system supports monitoring and correlation and does \textbf{not} predict tsunamis.}

\subsection{Geographic and Temporal Scope}
The project covers the Pacific basin, including Alaska and Hawaii, during 2010--2025, with earthquakes of magnitude $M \geq 6$.
NOAA HazEL regions 80--89 and a USGS rectangle (latitude $-60$ to $65$, longitude $100$ to $300$) provide geographic filters; the USGS historical sample covers 2010--2011 to meet the workshop's 200-record requirement.
The two real-time feeds are profiled in full, with Pacific events identified by coordinates; tsunami.gov covers the Pacific, Alaska, the US and the Caribbean.

\begin{tabularx}{\textwidth}{@{}lXX@{}}
\toprule
Source & Coverage & Fits the scope? \\
\midrule
NOAA HazEL & Global & Yes, regions 80--89 \\
USGS FDSN & Global & Yes, Pacific rectangle \\
USGS real-time feeds & Global & Yes, in-scope events identified by coordinates \\
tsunami.gov & Pacific, Alaska, US, Caribbean & Yes, in-scope events identified by coordinates \\
\bottomrule
\end{tabularx}

\subsection{Meaning of ``Current Situation''}
\todo{For example: earthquakes of the last 24 hours and 7 days, and the latest tsunami.gov message.}



% ===== 02_historical.tex =====
% Owner: TBD
\section{Historical Sources}
\label{sec:historical}

\subsection{NOAA/NCEI HazEL -- Global Historical Tsunami Database}
\begin{tabularx}{\textwidth}{@{}lX@{}}
\toprule
Owner & \todo{NOAA National Centers for Environmental Information (NCEI)} \\
URL & \url{https://www.ngdc.noaa.gov/hazel/view/swagger} \\
Access method & REST API, JSON \\
Authentication & None required \\
Rate limits & \todo{verify and cite} \\
License and attribution & \todo{verify and cite} \\
Fields used & \todo{list, e.g.\ year, month, day, latitude, longitude, causeCode, eventValidity, eqMagnitude} \\
Time coverage & \todo{from the catalog start to the present} \\
Update policy & \todo{} \\
\bottomrule
\end{tabularx}

\paragraph{Evidence.} \todo{Real response snippet (see Appendix~\ref{app:snippets}) and the file under \path{workshop_1/samples/}.}

\paragraph{Volume and partitioning.} \todo{Records in scope (2010--2025, regions 80--89), size, pagination of 200 records per page, and how the bulk load will be partitioned.}

\subsection{USGS FDSN Event Web Service}
\begin{tabularx}{\textwidth}{@{}lX@{}}
\toprule
Owner & U.S.\ Geological Survey, Earthquake Hazards Program \\
URL & \url{https://earthquake.usgs.gov/fdsnws/event/1/} \\
Access method & REST API, CSV (also GeoJSON and QuakeML) \\
Authentication & None required \\
Rate limits & \todo{maximum 20,000 events per request; cite documentation} \\
License and attribution & \todo{verify and cite} \\
Fields used & \todo{id, time, latitude, longitude, depth, mag, magType, status, updated} \\
Time coverage & \todo{} \\
Update policy & \todo{reviewed events are revised later; see Section~\ref{sec:incremental}} \\
\bottomrule
\end{tabularx}

\paragraph{Evidence.} \todo{}
\paragraph{Volume and partitioning.} \todo{Estimate for 2010--2025 and the decision to load one year at a time.}



% ===== 03_incremental.tex =====
% Owner: TBD
\section{Incremental Sources}
\label{sec:incremental}

\subsection{USGS Real-Time GeoJSON Feed}
\begin{tabularx}{\textwidth}{@{}lX@{}}
\toprule
URL & \url{https://earthquake.usgs.gov/earthquakes/feed/v1.0/geojson.php} \\
Feed used & Magnitude 4.5+, past 30 days \\
Format & GeoJSON \\
Update frequency & \todo{the server sets a cache of 60 seconds} \\
Change detection & \todo{\texttt{updated} field, stable \texttt{id}, Last-Modified header} \\
\bottomrule
\end{tabularx}

\subsection{tsunami.gov Warning Centers}
\begin{tabularx}{\textwidth}{@{}lX@{}}
\toprule
URL (NTWC) & \url{https://www.tsunami.gov/events/xml/PAAQAtom.xml} \\
URL (PTWC) & \url{https://www.tsunami.gov/events/xml/PHEBAtom.xml} \\
Format & Atom feed; each entry links to a CAP message \\
Content & \todo{one latest message per feed; most messages are information statements} \\
Change detection & \todo{entry id, \texttt{updated}, Last-Modified header} \\
\bottomrule
\end{tabularx}

\subsection{Empirical Update Behavior}
Two snapshots were taken on 7 November 2026 at 05:04 and 14:03 UTC (about 9 hours apart).

\begin{table}[H]
\centering
\begin{tabular}{lrr}
\toprule
 & USGS feed & tsunami.gov (NTWC, PTWC) \\
\midrule
First snapshot (records) & 714 & 1 per feed \\
Second snapshot (records) & 726 & 1 per feed \\
New & 22 & 0 \\
Modified & 25 & 0 \\
Removed & 10 & 0 \\
\bottomrule
\end{tabular}
\caption{Changes between the two snapshots.}
\end{table}

\todo{Explain: all 10 removed events left the rolling 30-day window (none was deleted); fields that changed in the modified events (time, place, magnitude, station counts); event ids stayed stable.}

\subsection{Revised or Deleted Records and Watermark Implications}
\todo{Watermark on \texttt{updated}, upsert on (source, id), content hash to separate real changes from timestamp-only updates, overlap window.}



% ===== 04_profiling.tex =====
% Owner: TBD
\section{Source Profiling and Data Quality}
\label{sec:profiling}

\subsection{Profiling Summary}
\todo{Rows, null percentages, ranges, duplicated identifiers, timestamp format and time zone, coordinate validity and units for each source. Reference the notebook \path{notebooks/01_source_profiling.ipynb}.}

\begin{tabularx}{\textwidth}{@{}lrrX@{}}
\toprule
Source & Records & Duplicated ids & Notes \\
\midrule
NOAA HazEL & \todo{} & \todo{} & \todo{} \\
USGS FDSN & \todo{} & \todo{} & \todo{} \\
USGS real-time & \todo{} & \todo{} & \todo{} \\
tsunami.gov & \todo{} & \todo{} & \todo{} \\
\bottomrule
\end{tabularx}

\subsection{Data-Quality Problems}
At least five, each with an example record.

\begin{tabularx}{\textwidth}{@{}clXX@{}}
\toprule
\# & Source & Problem & Example record \\
\midrule
1 & NOAA HazEL & Official warning status missing & \todo{} \\
2 & NOAA HazEL & Erroneous or doubtful events & \todo{} \\
3 & NOAA HazEL & Time of day or coordinates missing & \todo{} \\
4 & USGS FDSN & Columns that are always empty & \todo{} \\
5 & tsunami.gov & Location as free text; depth units differ & \todo{} \\
\bottomrule
\end{tabularx}

\subsection{Cross-Source Overlap}
At least five candidate pairs describing the same real-world event.

\begin{tabularx}{\textwidth}{@{}clllllX@{}}
\toprule
\# & Source A id & Source B id & $\Delta t$ & Distance & Magnitude & What differs \\
\midrule
1 & PAAQ-1-to0j1g & nc75462637 & $-1.3$ s & 8.7 km & 4.5 Ml / 4.96 mw & magnitude type, depth, location text \\
2 & \todo{} & \todo{} & \todo{} & \todo{} & \todo{} & \todo{} \\
3 & \todo{} & \todo{} & \todo{} & \todo{} & \todo{} & \todo{} \\
4 & \todo{} & \todo{} & \todo{} & \todo{} & \todo{} & \todo{} \\
5 & \todo{} & \todo{} & \todo{} & \todo{} & \todo{} & \todo{} \\
\bottomrule
\end{tabularx}



% ===== 05_mapping.tex =====
% Owner: TBD
\section{Canonical Mapping}
\label{sec:mapping}

\subsection{Field Mapping to the Canonical \texttt{event}}
\begin{tabularx}{\textwidth}{@{}lXXXX@{}}
\toprule
Canonical field & NOAA HazEL & USGS FDSN / feed & tsunami.gov (CAP) & Rule \\
\midrule
\texttt{event\_id} & \todo{} & \todo{} & \todo{} & \todo{} \\
\texttt{source} & \todo{} & \todo{} & \todo{} & \todo{} \\
\texttt{event\_type} & \todo{} & \todo{} & \todo{} & \todo{} \\
\texttt{timestamp} (UTC) & \todo{} & \todo{} & \todo{} & \todo{} \\
\texttt{latitude} & \todo{} & \todo{} & \todo{} & WGS84 \\
\texttt{longitude} & \todo{} & \todo{} & \todo{} & WGS84, $-180$ to $180$ \\
\texttt{severity} & \todo{} & \todo{} & \todo{} & \todo{value and unit/type} \\
\texttt{raw\_reference} & \todo{} & \todo{} & \todo{} & id or URL plus content hash \\
\bottomrule
\end{tabularx}

\subsection{Domain Tables}
\todo{Proposed tables (for example \texttt{earthquake}, \texttt{tsunami\_event}, \texttt{tsunami\_alert}) with their attributes.}

\subsection{Unit Conversions and Rules}
\todo{Timestamps (UTC), coordinates (WGS84), longitude normalization, depth units, magnitude types and severity scale.}

\subsection{Draft ER Diagram}
\begin{figure}[H]
\centering
\fbox{\parbox{0.8\textwidth}{\centering\vspace{3cm}\todo{insert ER diagram}\vspace{3cm}}}
% \includegraphics[width=\textwidth]{figures/er_diagram.png}
\caption{First draft of the entity-relationship diagram.}
\end{figure}



% ===== 06_contract.tex =====
% Owner: TBD
\section{Ingestion Contract Draft}
\label{sec:contract}

\begin{tabularx}{\textwidth}{@{}lXXXX@{}}
\toprule
Source & Historical strategy & Incremental strategy & Duplicate key & Overlap window \\
\midrule
NOAA HazEL & \todo{} & \todo{} & \todo{} & \todo{} \\
USGS FDSN & \todo{} & \todo{} & \todo{} & \todo{} \\
USGS real-time feed & -- & \todo{} & \todo{} & \todo{} \\
tsunami.gov & -- & \todo{} & \todo{} & \todo{} \\
\bottomrule
\end{tabularx}

\subsection{Correctness After a Failed Run}
\todo{The watermark advances only after a successful commit; idempotent writes; what happens to a partially loaded batch.}



% ===== 07_project.tex =====
% Owner: TBD
\section{Project Definition}
\label{sec:project}

\subsection{Project Overview}
\todo{Title, motivation, goal (data-engineering and analytical), scope, explicit exclusions and constraints.}

\subsection{Functional User}
\todo{Role, profile, decisions supported, and why the dashboard is useful.}

\subsection{User Stories}
\begin{tabularx}{\textwidth}{@{}clXc@{}}
\toprule
ID & Type & Story & MoSCoW \\
\midrule
US1 & Monitoring & As a \todo{user}, I want to \todo{action} so that \todo{benefit}. & \todo{} \\
US2 & Monitoring & \todo{} & \todo{} \\
US3 & Analysis & \todo{} & \todo{} \\
US4 & Analysis & \todo{} & \todo{} \\
US5 & Trust & \todo{} & \todo{} \\
US6 & Trust & \todo{} & \todo{} \\
\bottomrule
\end{tabularx}

\subsection{Indicator Proposal}
\todo{Name, inputs (magnitude, depth, location, historical tsunami association, official alert), rule, output scale, how it will be validated against historical events. State clearly that it is a monitoring aid and not a prediction, and list its limitations.}

\subsection{Risks and Mitigation}
\begin{tabularx}{\textwidth}{@{}lXX@{}}
\toprule
Risk & Impact & Mitigation \\
\midrule
Source outage & \todo{} & \todo{} \\
Rate limits & \todo{} & \todo{} \\
Schema change & \todo{} & \todo{} \\
License and attribution & \todo{} & \todo{} \\
Ambiguous identifiers & \todo{} & \todo{} \\
\bottomrule
\end{tabularx}

\subsection{Architecture Overview}
\begin{figure}[H]
\centering
\fbox{\parbox{0.8\textwidth}{\centering\vspace{3cm}\todo{insert architecture diagram}\vspace{3cm}}}
% \includegraphics[width=\textwidth]{figures/architecture.png}
\caption{High-level architecture adapted to tsunami monitoring.}
\end{figure}



% ===== 08_conclusions.tex =====
% Owner: TBD
\section*{Conclusions and Open Questions}
\addcontentsline{toc}{section}{Conclusions and Open Questions}
\todo{Main decisions taken in this workshop and what remains open for Workshop 2.}



% ===== 09_references.tex =====
% Owner: TBD
\begin{thebibliography}{9}
\bibitem{hazel} NOAA National Centers for Environmental Information. \emph{Hazard Event Database (HazEL), Global Historical Tsunami Database}. \url{https://www.ngdc.noaa.gov/hazel/view/swagger}. Accessed \todo{date}.
\bibitem{usgsfdsn} U.S.\ Geological Survey. \emph{API Documentation -- Earthquake Catalog}. \url{https://earthquake.usgs.gov/fdsnws/event/1/}. Accessed \todo{date}.
\bibitem{usgsfeed} U.S.\ Geological Survey. \emph{Real-time Notifications, Feeds, and Web Services (GeoJSON summary feed)}. \url{https://earthquake.usgs.gov/earthquakes/feed/v1.0/geojson.php}. Accessed \todo{date}.
\bibitem{tsunamigov} NOAA National Weather Service, Tsunami Warning Centers. \emph{Tsunami event feeds}. \url{https://www.tsunami.gov/}. Accessed \todo{date}.
\end{thebibliography}



\appendix

% ===== 10_appendix.tex =====
% Owner: TBD
\section{Real Response Snippets}
\label{app:snippets}
\todo{Short excerpts of the unmodified responses, one per source, with the file name in \path{workshop_1/samples/}.}

\begin{lstlisting}
(paste a short excerpt here)
\end{lstlisting}



\end{document}
